\documentclass[12pt, a4paper]{article}
\usepackage[margin=2.5cm]{geometry}
\usepackage{graphicx}
\usepackage{setspace}
\usepackage{parskip}
\usepackage{titlesec}
\usepackage{hyperref}
\usepackage{apacite}
\usepackage{natbib}
\usepackage{times}
\usepackage{microtype}
\usepackage[none]{hyphenat}
\tolerance=1
\emergencystretch=\maxdimen
\hyphenpenalty=10000
\hbadness=10000

\onehalfspacing

\title{\textbf{Blockmediary: Competitor Analysis}}
\author{}
\date{BEEM063 FinTech Hackathon \\ University of Exeter}

\begin{document}

\maketitle
\newpage

%---------------------------------------------------------------
\section*{Introduction}
%---------------------------------------------------------------

This assessment compares Blockmediary against the closest platforms operating in
programmable escrow and cross border trade finance, and sets that comparison against the
macro environment Blockmediary operates within. It was compiled to support the Main Group
Video pitch and the Individual Report, and feeds the competitive analysis and market
strategy narrative used across both. Figures throughout are drawn from public research
current as of August 2026 and should be treated as time sensitive, since funding and
licensing status in this sector changes quickly.

%---------------------------------------------------------------
\section*{1. Market Context}
%---------------------------------------------------------------

Cross border trade finance for small and medium sized enterprises is a large and
underserved market. The addressable market is projected to grow from approximately fifty
point nine billion United States dollars in 2026 to approximately eighty nine point five
billion dollars by 2034, a compound annual growth rate of seven point three per cent
\citep{marketsme2026}, and the underlying reason the market remains underserved is
structural: traditional Letters of Credit were built for large corporates, not small and
medium sized enterprises \citep{weforum2025}. Two structurally different approaches
currently compete for this market. The first is the bank consortium blockchain network,
exemplified by Komgo, Contour, we.trade, and the Marco Polo Network, each of which
digitises the existing letter of credit or bank workflow while keeping banks in the loop.
The second is the fintech escrow platform, exemplified by Tazapay and Truzo, which removes
the letter of credit entirely but keeps a centralised, custodial intermediary, namely the
fintech itself, together with manual document review. Blockmediary's wager is that neither
camp has combined full bank removal, non custodial on chain settlement, and a deterministic
document verification engine, with artificial intelligence extraction and a human review
console on the near term roadmap, in a single product, and that this combination is the
open lane in the market.

%---------------------------------------------------------------
\section*{2. Head to Head Comparison}
%---------------------------------------------------------------

Across the six platforms considered in this assessment, meaning Blockmediary and its five
closest comparators, several dimensions separate Blockmediary from the field. On the
question of whether a bank remains anywhere in the payment flow, Blockmediary removes
banks entirely, whereas Tazapay and Truzo continue to settle over bank or electronic money
institution rails, XREX is a licensed intermediary that only partially removes bank
dependency, and Komgo and Contour both retain a bank issued letter of credit at the centre
of the transaction, with the underlying technology only making that existing process
faster. On custody, Blockmediary is non custodial by design, since funds sit in a smart
contract rather than in an account Blockmediary controls, while Tazapay, Truzo, and XREX
are all custodial in the ordinary sense, Komgo is software only and never touches funds,
and Contour's funds moved through the participating banks. On settlement rail,
Blockmediary settles on chain using USDC on the Base network, whereas Tazapay uses SWIFT
and local payment rails with a recent stablecoin add on, Truzo settles in fiat through
Currencycloud, XREX settles across a multi asset custodial stack including Bitcoin,
Ethereum, Solana, and stablecoins, Komgo settles through traditional banking rails, and
Contour used R3 Corda purely for document exchange while settlement itself remained a bank
function. On document and trust verification, Blockmediary uses a deterministic rules
engine, with artificial intelligence extraction and a human review console on the roadmap,
whereas Tazapay relies on manual back office review, Truzo and
XREX rely on manual know your customer and compliance checks, with XREX supplementing this
with automated wallet risk scoring, and Komgo and Contour both rely on conventional bank
underwriting. Blockmediary is also the only one of the six with open source smart
contracts. On target customer, Blockmediary serves small and medium sized enterprise
buyers and sellers directly, Tazapay is embedded into business to business marketplaces,
Truzo is corridor focused on the United Kingdom and South Africa, XREX targets emerging
market small and medium sized enterprises with an Asia Pacific focus, and Komgo and
Contour both serve banks and large commodity or corporate clients rather than small and
medium sized enterprises directly. On fees, Blockmediary prices by document verification
tier rather than a single flat rate: zero point eight per cent for Tier A, verified through
electronic bills of lading, one point five per cent for Tier B, verified through carrier
application programming interfaces, and three per cent for Tier C, verified manually on
paper, validated in the deal value research report and financial model prepared for the
team in August 2026 \citep{dealvaluereport2026}. This compares against Tazapay's zero
point eight to three point eight per cent plus fixed charges, Truzo's zero point two two
to three point one per cent tiered structure plus foreign exchange charges, XREX's sub one
per cent transfer fee alongside a flat ten United States dollar Tether equivalent escrow
charge, Komgo's undisclosed enterprise licensing fees, and the roughly one to three per
cent fee typical of a bank issued letter of credit that Contour's model still carried. The
blended take rate across Blockmediary's tier mix starts near two per cent in Year one and
falls toward one per cent by Year five as adoption of the fastest, cheapest tier grows from
fifteen per cent to seventy per cent of volume. On funding and status, Tazapay has raised approximately
fifty seven point nine million dollars \citep{tazapayseriesa2023, tazapayseriesb2026} and
remains active and expanding, Truzo has raised a comparatively thin approximately six
hundred and forty thousand dollars with its last disclosed round in 2020 and remains
active but narrow in its corridor \citep{truzouklaunch2023}, XREX has raised approximately
forty two point seven million dollars including an eighteen point seven five million
dollar strategic investment from Tether and continues to grow \citep{xrextether2024}, Komgo
has raised over twenty six million Swiss francs from its bank and corporate shareholders
and remains active with over four hundred corporate and institutional clients
\citep{komgonewsroom2025}, and Contour raised approximately eighteen point three million
dollars before shutting down in November 2023 \citep{contourshutdown2023}. Blockmediary
itself remains in development, without the funding or licensing history the incumbents
have accumulated, which is a genuine gap this assessment does not minimise.

%---------------------------------------------------------------
\section*{3. PESTLE Analysis: Blockmediary's Operating Environment}
%---------------------------------------------------------------

A PESTLE analysis examines the macro environment a business operates within, namely
political, economic, social, technological, legal, and environmental factors, and is
conducted once for Blockmediary rather than once per competitor, since these factors are
largely external and shared across the market rather than being company specific in the
way a strengths and weaknesses analysis would be. The one axis that genuinely varies by
company, being regulatory regime, is folded into the political and legal discussion below
as a direct comparison rather than repeated five times over.

\subsection*{Political}

Favourable trade agreements underpin Blockmediary's chosen beachhead corridor: the United
Arab Emirates to India Comprehensive Economic Partnership Agreement has been active since
2022 and the United Arab Emirates to Turkiye equivalent has been in force since September
2023 \citep{uaeindiacepa2022, uaeturkiyecepa2023}, giving the chosen ten to fifty thousand
dollar textiles and garments corridor political tailwind through high trade volume and low
tariff friction. Corridor exclusions in this market are politically driven rather than
purely commercial. China was struck from the corridor mix because of the People's Bank of
China's 2021 ban on crypto asset speculation, which prevents Chinese small and medium
sized enterprise sellers from legally holding or receiving USDC \citep{pbocban2021}, and
Turkiye's stablecoin payments ban under Central Bank Regulation No. 2021/14 is a similar
political constraint, which is why Turkiye is excluded as a target market even though the
United Arab Emirates to Turkiye trade corridor itself remains in scope
\citep{cbrt202114}. Regulators in the target markets are, on the whole, actively building
on ramps for firms like Blockmediary rather than simply gatekeeping: the Financial
Services Regulatory Authority's Innovation Programme, known as RegLab, and the Dubai
International Financial Centre's FinTech Hive are both politically motivated to attract
crypto and fintech innovation to the United Arab Emirates free zones, which is why those
free zones, rather than the United Arab Emirates mainland or Dubai's Virtual Assets
Regulatory Authority regime, are the planned regulatory home. Postures diverge across
target jurisdictions in a way that matters for market entry sequencing: United Kingdom and
European Union regulators are still mid transition toward comprehensive crypto regimes,
with timelines that continue to slip, whereas United Arab Emirates free zone regulators
are more overtly courting fintech entrants, which favours the United Arab Emirates first
market entry sequencing already chosen. Sanctions regimes are a live political variable on
every individual deal rather than a one off consideration, since lists maintained by the
Office of Foreign Assets Control, the United Nations Security Council, His Majesty's
Treasury, and the United Arab Emirates can change between deal creation and fund release.
On regulatory regime specifically, the closest comparators sit under materially different
postures from Blockmediary's own: Truzo operates under United Kingdom Financial Conduct
Authority and South African Financial Sector Conduct Authority oversight, XREX operates
under Singapore and Taiwan licensing, and Komgo operates under European Union and Swiss
banking consortium oversight, and none of the five researched competitors is pursuing the
free zone first strategy Blockmediary is, which is a distinct, and arguably underexploited,
political positioning choice.

\subsection*{Economic}

The core economic opportunity here is a market gap rather than simply a large market size.
The ten to fifty thousand dollar deal size Blockmediary targets is what the team refers to
internally as the letter of credit dead zone: too small for banks to serve economically
under a traditional letter of credit, and too large to be handled on trust alone. Every
researched competitor, being Tazapay, Truzo, XREX, and Komgo, is priced or structured for
a larger or otherwise different segment, and none targets this exact dead zone directly.
Currency instability is a direct economic driver of adoption in the beachhead corridor:
Turkish lira inflation has already pushed Turkish small and medium sized enterprises
toward stablecoin native behaviour, which means less user education is required to make
counterparties comfortable with USDC settlement, an economic tailwind specific to the
United Arab Emirates to Turkiye leg of the corridor. Fee economics are the headline
competitive lever. Traditional letter of credit fees run one to three per cent, and the
researched competitors cluster in a similar range, Tazapay at zero point eight to three
point eight per cent, Truzo at zero point two two to three point one per cent, and XREX
below one per cent on transfers though its escrow product separately charges a flat ten
United States dollar Tether equivalent fee. Blockmediary prices by verification tier at
zero point eight, one point five and three per cent, validated in the deal value research
report and financial model \citep{dealvaluereport2026}; removing the bank still removes
the largest cost line on the fastest, cheapest tier, undercutting the letter of credit
range at the low end, and the blended take rate across the tier mix falls from near two
per cent in Year one toward one per cent by Year five as adoption shifts toward that
fastest tier, a structural cost advantage on the deals that matter most rather than a
flat, aspirational claim. The underlying pain point is a small and medium sized enterprise cash
flow problem: the roughly ninety day payment delay pattern typical of cross border trade,
in which a seller ships and then waits for letter of credit processing and payment,
directly damages small and medium sized enterprise working capital, so Blockmediary's near
instant, rules based release changes the economics of running a
trading business, not merely the cost of a single transaction. A genuine headwind should
also be acknowledged rather than glossed over: a higher global interest rate environment
increases the opportunity cost to a buyer of locking capital in escrow before goods arrive,
which is a real economic cost of any prefunded escrow model, Blockmediary included.

\subsection*{Social}

Trust is a social problem here as much as a technical one. Small and medium sized
enterprise owners in the target corridors have historically low trust both in crypto,
often perceived as scam adjacent, and in an unfamiliar counterparty in a new country.
Blockmediary is competing on the same social trust gap that Truzo's anti fraud branding
and Tazapay's transparent fee messaging are also trying to close, but Blockmediary's
approach is automated, on chain proof rather than brand reassurance. Diaspora networks
function as a genuine social distribution channel rather than simply a demographic fact:
the large Indian diaspora resident in the United Arab Emirates underpins the warm
introduction go to market motion in the beachhead corridor, since social trust travels
faster through existing community networks than through cold outreach or advertising.
Generational and digital literacy shifts favour adoption, since younger small and medium
sized enterprise owners and operations staff are more comfortable transacting in
stablecoins and using app based workflows than the generation that built trust
relationships with relationship bankers, which is what makes a non bank, app first
alternative to a letter of credit socially plausible now in a way it was not a decade ago.
Founder led business development depends on social capital with freight forwarders, who
see the small and medium sized enterprise trust gap problem earlier than banks do and have
no commercial reason not to refer a deal, a relationship dynamic the go to market plan
deliberately leans on in preference to paid acquisition channels.

\subsection*{Technological}

Two specific technology maturations make Blockmediary possible now in a way it would not
have been a few years ago. First, low cost, fast finality layer two blockchains such as
Base make micro value on chain escrow economically viable, where the equivalent
transaction would have been cost prohibitive on the Ethereum base layer. Second, large
language model based document understanding now makes automated verification of
commercial invoices, bills of lading, and packing lists against a rules engine accurate
enough to meaningfully reduce, though not eliminate, human review, in a way earlier
optical character recognition only approaches could not, since they could not handle the
judgement calls that Uniform Customs and Practice for Documentary Credits style document
checking requires. Blockmediary's delivered prototype currently uses the deterministic
rules engine on its own; this large language model extraction layer is a near term
roadmap item rather than delivered functionality. Legal recognition of electronic trade documents is catching up to the
underlying technology: the United Kingdom's Electronic Trade Documents Act 2023 and the
United Nations Commission on International Trade Law's Model Law on Electronic
Transferable Records, already adopted in the Dubai International Financial Centre,
Singapore, and Bahrain, are what will eventually let Blockmediary move beyond scanned
portable document format bills of lading to fully electronic, title conferring electronic
bills of lading \citep{etda2023, uncitralmletr}, though this remains a near term roadmap
item gated by the convergence of law and technology rather than by engineering alone. It
is worth noting explicitly that the prior generation of blockchain trade finance
technology failed for organisational reasons rather than technical ones: Komgo, Contour,
we.trade, and the Marco Polo Network all had working distributed ledger technology,
whether Corda, Quorum, or Hyperledger Fabric, and instead failed on adoption economics and
consortium governance, discussed further in section four. The lesson for Blockmediary's
own technology strategy is that technical sophistication is necessary but not sufficient,
and that the architecture must also avoid requiring closed network infrastructure in which
both sides of a transaction need to be onboarded to the same platform. A further point
worth stating as a risk rather than only an enabler is that the releaser private key
represents a single point of failure, identified in Blockmediary's own risk register as
the platform's top security risk \citep{internallegalrisk2026}; production deployment
requires a hardware security module or threshold signature scheme rather than a hot
wallet, which remains a live technological constraint on scaling past the minimum viable
product. Finally, cloud and vendor infrastructure, including know your business providers,
large language model application programming interfaces, and blockchain remote procedure
call nodes, is now mature and accessible enough for a small team to assemble a compliant
technology stack quickly, which lowers the barrier to entry for Blockmediary itself but
equally lowers it for any future copycat competitor.

\subsection*{Legal}

The regulatory picture facing Blockmediary is a genuine patchwork rather than a single
hurdle. The United Kingdom imposes stablecoin custody rules and cryptoasset service
provider licensing requirements, the European Union imposes Markets in Crypto Assets
Regulation authorisation, mandatory from July 2026, alongside Digital Operational
Resilience Act obligations, and the United Arab Emirates imposes the Central Bank's
Payment Token Services Regulation on the mainland and the Virtual Assets Regulatory
Authority's Rulebook in Dubai, against the comparatively lighter touch free zone
frameworks in the Dubai International Financial Centre and the Abu Dhabi Global Market;
each jurisdiction imposes a different licensing threshold, and the full picture is set out
in Blockmediary's own thirty one item risk register \citep{internallegalrisk2026}. Anti
money laundering, know your customer, and sanctions compliance form a mandatory gate on
every individual deal rather than a one time cost, since customer due diligence and
screening obligations must be satisfied before funding and sanctions lists must be
re-checked at both funding and release, which is an ongoing operational and legal cost
baked directly into Blockmediary's unit economics. Smart contract enforceability remains
legally unsettled in the United Kingdom, and Blockmediary mitigates this with a separate,
signed Trade Escrow Agreement that governs legal rights, with the smart contract itself
serving only as the execution layer; none of the researched blockchain trade finance
competitors that failed, being Contour, we.trade, and the Marco Polo Network, solved this
question any differently, which suggests it is an industry wide open question rather than
a Blockmediary specific gap. Data protection law directly shapes the product's
architecture: the General Data Protection Regulation's right to erasure is fundamentally
incompatible with writing personal data to an immutable public chain, which is why the
confirmed architecture keeps all personal data off chain and commits only hashes and
wallet addresses on chain, a legal constraint that became a core design decision rather
than an afterthought. Legal recognition gaps directly limit product scope today: without
fully integrated Electronic Trade Documents Act or Model Law on Electronic Transferable
Records qualifying electronic bill of lading systems, Blockmediary's minimum viable product
accepts scanned portable document format documents only, which sets a legal maturity
ceiling on how far the paperless pitch can go until electronic bill of lading
infrastructure and law mature together. Finally, legal exclusions define the addressable
market as much as economics do: Turkiye and mainland China are excluded primarily because
their respective stablecoin payment bans make Blockmediary's core product illegal there,
not because the underlying market opportunity is unattractive.

\subsection*{Environmental}

Chain choice carries a genuine environmental, social, and governance angle worth using in
the pitch. Base runs on Ethereum's proof of stake consensus mechanism, so Blockmediary's on
chain settlement layer carries a materially smaller energy footprint than proof of work
alternatives such as Bitcoin, which is a minor but genuine differentiator for judges or
investors screening for environmentally conscious infrastructure choices. Digitising trade
documentation also reduces physical logistics overhead: replacing courier driven paper
document exchange, a real cost and emissions source in traditional letter of credit and
documentary trade, with digital upload and verification is a modest environmental
positive, secondary to the cost and speed argument but usable as supporting colour.
Climate linked supply chain disruption is an indirect but real risk factor: extreme
weather affecting shipping routes and ports, for example Red Sea or Suez disruptions or
storm driven port congestion, affects the underlying trade volumes and shipment timelines
that Blockmediary's document release logic depends upon, a factor outside Blockmediary's
control but worth a brief acknowledgement in any risk discussion. Environmental regulation
is, at present, the lightest of the six factors for Blockmediary at minimum viable product
stage: unlike the political, legal, and economic factors above, carbon reporting mandates
and green trade finance requirements are not yet a binding constraint on the beachhead
corridor, so this axis is best treated as a forward looking consideration rather than a
current risk.

%---------------------------------------------------------------
\section*{4. Company Profiles}
%---------------------------------------------------------------

\subsection*{4.1 Tazapay}

Tazapay is headquartered in Singapore, was founded in 2020, and has raised approximately
fifty seven point nine million dollars, including a Series A led by Sequoia Southeast Asia
and an extended Series B of approximately thirty six million dollars
\citep{tazapayseriesa2023, tazapayseriesb2026}. It began as a payment gateway and now runs
a flagship escrow as a service product embedded by application programming interface into
business to business marketplaces spanning freight, gaming, travel, and small and medium
sized enterprise import and export \citep{tazapayeaas2021}. It accepts payments in over
one hundred and seventy three countries and pays out to over seventy markets. It has
recently added stablecoin on and off ramps, but its core rails remain traditional banking
over SWIFT and local networks, it is licensed by the Monetary Authority of Singapore, and
it is fully custodial. Document and trust verification is manual: Tazapay staff check
shipping documents before releasing funds, with a turnaround of roughly forty eight hours,
and fees run zero point eight to three point eight per cent plus fixed charges. Its
Trustpilot rating sits at approximately two out of five across four hundred and ninety
five reviews, with recurring complaints about refund friction, arbitrary account
suspensions, multi day payment freezes during fraud review, and, most damaging for an
escrow brand specifically, allegations that funds were sometimes released to suppliers
before goods had actually arrived \citep{tazapaytrustpilot}. Tazapay is, in substance, the
clearest same problem, old architecture foil to Blockmediary: centralised, custodial, and
human reviewed. Blockmediary targets precisely its weakest points, turning opacity and
slow, arbitrary release into instant, rules based, on chain release.

\subsection*{4.2 Truzo}

Truzo is headquartered across Johannesburg and London, was founded in 2017, and has raised
a comparatively thin approximately six hundred and forty thousand dollars, with its last
disclosed round reported in 2020. It is a pure escrow as a service wallet, originally
built for South African peer to peer and business anti fraud use cases, and expanded to
the United Kingdom in 2023, becoming Financial Conduct Authority approved specifically to
serve the United Kingdom to South Africa trade corridor, worth roughly ten billion pounds
sterling a year \citep{truzouklaunch2023}. Fees are tiered by currency and transaction
size, running zero point two two to three point one per cent, with an additional twenty
five pound intermediary bank surcharge for non United Kingdom buyers \citep{truzofees}. It
runs on traditional fiat rails through Currencycloud, with no blockchain and no
stablecoins, and its verification is manual know your customer and anti money laundering
compliance rather than document driven trust. Truzo shows a thin funding and traction
history for a nine year old company, with its last raise reported in 2020 and only around
six Trustpilot reviews \citep{truzotrustpilot}; its geographic corridor is narrow, and its
fees at low transaction tiers, up to three point one per cent, are comparable to or worse
than the letter of credit fees Blockmediary is trying to undercut. Truzo is, in substance,
a regulated fintech escrow wallet rather than a trade finance product, with no document
verification layer at all, and Blockmediary's pitch of removing banks entirely and pricing
by verification tier is a direct rebuttal to Truzo's still bank routed, flat percentage
fee model. Truzo's absolute floor, zero point two two per cent, only available above
roughly ten million pounds sterling equivalent, can undercut Blockmediary's Tier A rate,
but Blockmediary's zero point eight, one point five and three per cent structure beats
Truzo's pricing at the deal sizes Truzo actually serves, and without the twenty five pound
intermediary bank surcharge Truzo adds for non United Kingdom buyers.

\subsection*{4.3 XREX}

XREX is headquartered in Taipei, was founded in 2018, and has raised approximately forty
two point seven million dollars, including an eighteen point seven five million dollar
strategic investment from Tether in 2024 \citep{xrextether2024}. It is a full blockchain
enabled financial institution, offering crypto fiat exchange, custody, cross border
payments, and a business to business escrow product called BitCheck, which holds crypto or
fiat until both parties confirm fulfilment, for a fee of roughly ten United States dollar
Tether equivalent \citep{xrexbitcheck}. Its XREX Pay product offers sub one per cent cross
border transfers using USDT and USDC settlement, explicitly targeting emerging market
small and medium sized enterprises facing United States dollar liquidity shortages. It has
a strong Asia Pacific and emerging market focus, and leans on heavy regulatory licensing,
including a Singapore Major Payment Institution licence and Taiwan Virtual Asset Service
Provider registration, as its trust signal. XREX is fundamentally custodial, holding funds
and keys, the opposite of Blockmediary's non custodial model, and its verification
combines manual know your customer checks with automated wallet risk scoring rather than
document based trust, with no evidence of invoice or bill of lading verification of any
kind, and no open source contracts. XREX is the closest of the five comparators in ethos,
being stablecoin first with an emerging market small and medium sized enterprise focus,
but it is architecturally opposite to Blockmediary on custody and trust model: Blockmediary
runs non custodial smart contracts with deterministic, rules based document verification,
with artificial intelligence extraction and a human review console on the roadmap, where
XREX runs a custodial, licensed intermediary with know your customer and wallet risk
scoring.

\subsection*{4.4 Komgo}

Komgo is headquartered in Geneva, was founded in 2018, and has raised over twenty six
million Swiss francs from a group of twenty to twenty five bank and corporate
shareholders, including ABN AMRO, BNP Paribas, Citi, HSBC, Shell, and TotalEnergies
\citep{komgoabout}. It is not a payments or escrow company at all: Komgo sells enterprise
software as a service to banks and commodity traders to digitise letters of credit,
guarantees, know your customer processes, and document exchange, and never touches funds
directly \citep{komgoconsensys}. Notably, Komgo abandoned its own blockchain, originally
built on Quorum, in favour of a centralised database for its core letter of credit
product, with its chief executive describing dropping blockchain as one of the best
decisions the company made; it remains fundamentally bank intermediated, since banks
continue to underwrite and hold funds while Komgo simply moves the paperwork faster. As of
2026 Komgo remains active and growing, with over four hundred corporate and institutional
clients, over sixty connected banks, and recent partnership wins including Standard
Chartered, Credit Agricole's Optimtrade, and the National Bank of Greece
\citep{komgonewsroom2025}. Komgo is the clearest old guard foil among the five comparators.
Its trajectory, retreating from blockchain back toward centralised infrastructure, is the
opposite of Blockmediary's own direction, and it never removes banks from the transaction
or serves small and medium sized enterprises directly, since such enterprises only reach
Komgo indirectly through their bank.

\subsection*{4.5 Contour: Shut Down 2023, a Case Study}

Contour was headquartered in Singapore, founded in 2020 as a spin out of the earlier
Voltron pilot on the R3 Corda platform, and raised approximately eighteen point three
million dollars from nine bank shareholders, including HSBC, Citi, Standard Chartered, and
BNP Paribas, together with R3, Bain, and CryptoBLK. It digitised letters of credit on R3
Corda, a permissioned bank distributed ledger, cutting letter of credit processing from
five to ten days down to under twenty four hours in pilot deployments. By 2023 it had
grown to nine investor banks plus twenty two member banks across seventeen or more
countries and had onboarded corporates including Posco International, yet actual usage was
only around sixty to seventy transactions per month network wide, far too low to cover the
cost of its enterprise infrastructure \citep{contourshutdown2023}. Contour shut down
between October and November 2023, and the reasons matter directly for Blockmediary's own
design choices. First, there was no lead investor: a coalition of competing banks with no
single accountable funder struggled to align on a fresh capital raise. Second, Contour
faced a closed network chicken and egg problem, since value only materialised once both
counterparty banks were onboarded to Corda, adoption stayed patchy, and most real trades
fell back to paper or SWIFT. Third, its revenue depended on bank goodwill rather than a
self serve customer base, which proved unsustainable at roughly sixty to seventy
transactions a month. Fourth, its permissioned infrastructure, priced per Corda node, was
expensive relative to cheap public chain rails. Contour's brand and intellectual property
were subsequently acquired by the fintech Xalts in February 2024, and then by XDC Ventures
in 2025 in an attempt to revive the brand \citep{contourxalts2024}, though the original
bank consortium model itself had already failed. The lessons apply directly to
Blockmediary's design: avoid consortium governance in favour of a single accountable
product rather than a coalition of competing banks; avoid any closed network onboarding
requirement, since Blockmediary's buyer and seller need only hold USDC on Base rather than
integrate with a counterparty bank; remove banks from the trust layer entirely, using a
deterministic rules engine and a smart contract, rather than merely routing trust through
issuing and advising banks faster; run a self serve fee model that does not depend on bank
subsidy or goodwill; and build on a public layer two network with standard USDC rails
rather than expensive permissioned enterprise distributed ledger technology. Contour's
failure is also not an isolated case. The equivalent European Union bank consortium
network we.trade, built on Hyperledger Fabric, entered liquidation in June 2022
\citep{wetradefuturum}; the Marco Polo Network, a roughly forty five bank consortium also
built on R3 Corda, filed for insolvency in February 2023 after a reported Bank of America
funding deal fell through \citep{marcopoloinsolvency2023}; and TradeLens, the IBM and
Maersk shipping logistics blockchain, was discontinued in November 2022, citing an
inability to achieve full global industry collaboration \citep{tradelens2022}. Taken
together, these four failures form a pattern rather than a series of one off events, and
that pattern, bank consortium ownership, a closed permissioned network, dependence on
every counterparty separately onboarding, and no single accountable driver of growth or
funding, is precisely the pattern Blockmediary's design deliberately avoids.

%---------------------------------------------------------------
\section*{5. Other Adjacent Players}
%---------------------------------------------------------------

A number of further platforms surfaced during this research without being core to the
five primary comparators, and are recorded briefly here rather than analysed in depth.
dltledgers is a blockchain platform for small and medium sized enterprise trade and supply
chain finance with direct overlap to Blockmediary's space, though it has not been
researched deeply for this assessment. Damisa is a blockchain and stablecoin based
business to business escrow product claiming an approximately eighty per cent cost
reduction relative to traditional trade finance, and is a close conceptual peer to
Blockmediary. Castler is an India based smart contract escrow infrastructure provider with
an enterprise compliance angle, functioning as a regional peer in the Indian market.
Vaultion is a claimed non custodial, multi chain escrow protocol reportedly supporting
Base among other networks, but it should be treated as unverified: the only source located
for it is a single startup directory listing, with no working website, source code
repository, named team, or funding data found anywhere else, and it should not be
presented as a substantiated competitor without that caveat attached
\citep{vaultionlisting}. WorldFirst, operating as WorldTrade, offers smart contract
business to business sourcing payment protection for small and medium sized enterprise
buyers and is an adjacent buyer protection product rather than a direct comparator.
EscrowLayer, Uniscrow, Escrowly, and Tscrow are generic crypto or stablecoin escrow
application programming interfaces that are not trade finance specific, functioning as
peer to peer or over the counter tooling rather than letter of credit replacements, and
are accordingly a lower priority for comparison purposes. Bridge, operated by Stripe,
BVNK, which is in the process of being acquired by Mastercard, Zero Hash, and Conduit are
stablecoin cross border payment infrastructure providers whose rails overlap with
Blockmediary's without offering an equivalent escrow or trust product. Ripple is a cross
border settlement network connected to over three hundred financial institutions, but it
is payments focused rather than escrow focused. Vakt, a sister project to Komgo, is a
blockchain post trade platform for physical commodities and is accordingly adjacent to
Komgo rather than to Blockmediary directly.

%---------------------------------------------------------------
\section*{6. Positioning Takeaways}
%---------------------------------------------------------------

Several conclusions follow from this comparison and should inform both the pitch deck and
the individual report. The market has already tried two failed or limited approaches:
bank consortiums, illustrated by Komgo's retreat from blockchain and by the Contour, we.
trade, Marco Polo, and TradeLens shutdowns, and centralised fintech escrow, illustrated by
Tazapay and Truzo's manual review, custodial risk, and fee levels that sit close to letter
of credit fees. Blockmediary is positioned as the synthesis neither camp has reached. The
absence of a lead investor and the wider problem of consortium governance is a proven
failure mode rather than a hypothetical risk, evidenced directly by Contour's own account
of its shutdown, and this is usable evidence that bank led blockchain trade finance
struggles structurally to move quickly, in a way a focused, non bank owned product does
not. Custody and verification are the two axes that most cleanly separate every researched
competitor from Blockmediary, since no researched competitor is both non custodial and
rules based document verified at the same time, which is why the head to head comparison in
section two is built around exactly those two dimensions. On fees specifically, Tazapay's
zero point eight to three point eight per cent and Truzo's zero point two two to three
point one per cent both land in or above the one to three per cent range typical of letter
of credit fees. Blockmediary's tiered zero point eight, one point five and three per cent
structure undercuts that range on Tier A and matches it on Tier B and Tier C, with a
blended take rate starting near two per cent in Year one and falling toward one per cent
by Year five as Tier A adoption grows, a useful side by side statistic for a pitch slide,
now grounded in the deal value research report and financial model rather than an
aspirational close to zero per cent claim. Finally, a caveat should
be stated explicitly rather than omitted from the individual report: several of the
platforms compared here, particularly Komgo and XREX, are larger, regulated, and well
funded incumbents, and Blockmediary's honest counter to that scale is speed of trust,
meaning rules based verification measured in minutes rather than bank underwriting
measured in days, and cost, meaning the absence of a bank margin, rather than a claim to
greater scale or regulatory maturity. Artificial intelligence extraction and a human
review console remain on the roadmap rather than delivered, which is itself a gap worth
acknowledging directly in the reflection section of the individual report, alongside the
scale and regulatory maturity points.

\newpage
\bibliographystyle{apacite}
\bibliography{Competitor}

\end{document}
